\documentclass[conference]{IEEEtran}
\usepackage[utf8]{inputenc}
\usepackage[T1]{fontenc}
\usepackage{cite}
\usepackage{amsmath,amssymb}
\usepackage{graphicx}
\usepackage{booktabs}
\usepackage{url}

\title{Evaluating a Deterministic Validator Plus Single‑Iteration Repair Pipeline for LLM‑Generated Network Configurations}

\author{
\IEEEauthorblockN{Autonomous AI Researcher}
\IEEEauthorblockA{CNSM 2026 Agentic AI Researcher Track}
}

\begin{document}

\maketitle

\begin{abstract}
We preregistered and executed a paired confirmatory experiment (N=300 paired intent\ensuremath{\rightarrow}configuration tasks; preregistration id cand-2026-001-A-prereg-v1, preregistration SHA-256 archived) to evaluate whether a minimal guarded pipeline---single‑shot LLM generation (declared model: gpt-5-mini) followed by a frozen deterministic configuration validator and at most one automated repair iteration---reduces first‑pass misconfiguration relative to plain single‑shot generation. Execution completed 600 episodes and used 301 model calls. Observed outcomes: baseline = 299/300 valid (99.667\%); guarded = 300/300 valid (100.0\%); paired counts n11=299, n10=1, n01=0, n00=0. Exact McNemar two‑sided p = 1.00; absolute paired difference = 0.0033333; paired bootstrap 95\% CI = [0.00, 0.01] (10,000 resamples, seed 1729). The single permitted automated repair corrected the sole invalid baseline output (repair success 1/1; Clopper--Pearson 95\% CI = [0.025, 1.0]). The preregistered planning assumption used for power calculations (planned baseline pass rate = 0.40) was contradicted by the observed near‑ceiling baseline, removing statistical headroom for the preregistered >=50\% relative reduction target. We report preregistered design choices, execution accounting, confirmatory statistics, a post‑hoc inferential sensitivity bound tied to the observed contingency counts, and artifact‑grounded recommendations for future NetOps confirmatory evaluations. All run artifacts and analysis outputs are archived in the execution manifest referenced in the Disclosure Statement.
\end{abstract}

\section{Introduction}
Translating operator intent to device configuration is a core NetOps task. Large language models (LLMs) have been proposed as intent\ensuremath{\rightarrow}configuration translators that can accelerate operations, but probabilistic generation risks misordering, omission, and unsafe commands. Recent literature advocates hybrid patterns that pair probabilistic generators with deterministic validators/simulators and bounded repair loops as operational floor‑safety nets. However, rigorous preregistered evaluations that quantify how tightly bounded guarded pipelines change first‑pass misconfiguration rates under controlled sampling semantics are limited.

We implement and preregister a narrowly scoped paired experiment to evaluate the operational effect of a minimal guarded pipeline composed of: (1) a single shared LLM candidate per task (single‑shot generation), (2) a frozen deterministic validator performing canonicalization and rule/dependency checks, and (3) at most one automated repair attempt when the validator reports violations. The core research question is: Can this minimal guarded pipeline reduce first‑pass misconfiguration rate relative to plain single‑shot generation under a controlled paired design? Our contributions are (i) a preregistered paired experiment with deterministic seed generation and archived per‑task artifacts, (ii) confirmatory statistics derived from those artifacts, (iii) a compact post‑hoc sensitivity quantification tied to the observed contingency counts, and (iv) artifact‑grounded recommendations for future NetOps confirmatory evaluations.

\section{Related Work}
Recent work on LLMs for NetOps and intent‑driven configuration has produced numerous system demonstrations and surveys documenting both promise and concrete failure modes. Empirical demonstrations of intent\ensuremath{\rightarrow}configuration translation report that LLMs can synthesize plausible device configuration templates in lab and testbed settings, yet they frequently surface operationally relevant errors such as command misordering, omission of required safety steps, and overconfident assertions that lack mechanical verification (example system evaluation and discussion: Nguyen et al. 2024 [10.1002/nem.2313]). Complementary survey and review articles consolidate these observations and emphasize the need for grounded verification and governance when applying LLMs to critical network tasks [NIST guidance and surveys, 10.6028/nist.ai.100-2e2023].

Parallel technical threads motivate guarded architectures. Work in adjacent domains (clinical QA, code repair, and general RAG/verification research) shows measurable reductions in factual errors when probabilistic generators are paired with retrieval, deterministic checks, or repair loops; these findings have been ported by several NetOps proposals arguing for generate\ensuremath{\rightarrow}verify\ensuremath{\rightarrow}repair pipelines rather than blind single‑shot deployment [e.g., comparative/procedural studies and framework discussions in the literature pool, 10.36227/techrxiv.176784505.56675782/v1; 10.2139/ssrn.6608518]. Architectural proposals for intent‑aware automated remediation and multi‑agent NetOps orchestration similarly argue that deterministic validators and bounded repair/backout policies are central to safe agentic NetOps workflows (see recent agentic NetOps proposals and intent‑translation demonstrations, e.g., 10.36227/techrxiv.177004277.71795055/v1).

Deterministic network verification provides concrete algorithmic foundations for validators used in guarded pipelines. Prior work on specification‑based and header‑space analyses demonstrates tractable static checks for reachability, access control, and update‑safety properties without requiring live device execution; such algorithms can be integrated to canonicalize generated artifacts into normalized ASTs and to flag violations that trigger repair logic [Beckett et al., 2017 foundation for incremental/specification checks, 10.1145/3098822.3098834]. These deterministic checks differ fundamentally from empirical or emulation tests: they provide provable, rule‑based pass/fail decisions for a defined specification, but they do not by themselves measure runtime control‑plane or packet‑level effects unless combined with higher‑fidelity simulation or emulation.

Limitations in the prior literature motivate the present work. Many demonstrations and proposals are system‑level or exploratory rather than preregistered confirmatory experiments; reported improvements from RAG, verifier‑assisted generation, or agentic architectures often lack paired within‑task statistical controls that isolate the incremental effect of a validator+repair floor on the same generated candidate. Moreover, published comparisons rarely quantify how sampling semantics (shared candidate versus independent draws) and repair budgets affect the number of informative discordant pairs available to paired binary tests. Surveys and architectural reviews therefore call for reproducible evaluation harnesses that combine deterministic validators with controlled sampling semantics and explicit inferential plans [10.6028/nist.ai.100-2e2023].

Our study addresses a tightly scoped portion of this gap by executing a preregistered paired experiment that (a) fixes shared\_initial\_candidate semantics to measure remediation capacity on a single generated plan, (b) uses a frozen deterministic validator to define a binary success criterion tied to intent constraints, and (c) limits repair budget to a single automated iteration. This design directly tests the operational claim that a minimal validator+repair floor reduces first‑pass misconfiguration, and it surfaces the inferential consequences of high baseline pass rates and shared‑candidate pairing for confirmatory McNemar‑style testing.

\section{Methodology}
Preregistration and hypotheses: The experiment was preregistered (cand-2026-001-A-prereg-v1) prior to observing outcomes. Confirmatory hypotheses were: H1 --- the guarded pipeline reduces first‑pass misconfiguration rate by >=50\% (relative) versus single‑shot generation; H2 --- one or fewer automated repair iterations recover >=75\% of initially invalid configurations. The preregistration explicitly used a planning assumption of planned baseline pass rate = 0.40 for power calculations; this numeric assumption is central to the interpretation of the confirmatory outcome and is documented in the preregistration record (see Disclosure).

Design and sampling semantics: We adopted a paired within‑task design with 300 deterministically generated seeds from a frozen task generator. For each seed the same sampled LLM candidate was reused in both arms (shared\_initial\_candidate semantics): baseline (single‑shot generation) and guarded (identical candidate validated and---if invalid---subjected to up to one automated repair attempt). Planned episodes = 600 (two episodes per seed), initial\_generation\_calls\_per\_task = 1, maximum\_repair\_calls\_per\_task = 1. The shared‑candidate pairing was preregistered to isolate the incremental effect of validation and repair on a fixed generated plan; this choice increases within‑pair correlation and reduces the number of informative discordant pairs available to McNemar inference compared with an independent‑draws design.

Model, validator, and scoring: The declared model in the execution contract was gpt-5-mini. The frozen deterministic validator/repair adapter (validator id deterministic\_netops\_validator\_v5) implements rule and dependency checks, canonicalizes configurations to normalized ASTs, and returns binary pass/fail scoring plus violation codes used to trigger and steer the preregistered automated repair adapter. The scoring rule is full intent‑constraint validation: outputs are scored as success only when required command sequences and workflow constraints are satisfied by the validator's canonicalization and rule checks. The preregistered missingness plan allowed one retry for provider failures; primary analysis used the 'complete\_pair\_primary' treatment which drops pairs only if both calls fail. All validator decisions, normalized configurations, and repair traces were archived in the execution manifest and analysis bundle.

Power, estimand, and analysis plan: The primary estimand was the paired\_success\_rate\_difference\_guarded\_minus\_baseline. The preregistered primary test was the exact McNemar two‑sided test (\ensuremath{\alpha}=0.05) on paired pass/fail outcomes; we also preregistered paired bootstrap percentile 95\% CIs (10,000 resamples) for the absolute paired difference. The preregistration included a planning assumption baseline pass rate = 0.40 which informed the sample size and the >=50\% relative reduction target; this assumption is reported here to make the planning calculus transparent. Secondary analyses included exact binomial intervals for repair success, stratified checks, and sensitivity analyses treating failed model calls as failures; all analysis code and per‑task artifacts are archived.

\section{Results}
Execution accounting and artifact provenance (archived): The run completed the planned 600 episodes (300 complete pairs) without provider or infrastructure failures under the preregistered retry policy. Total model\_calls\_used recorded in the execution manifest = 301, reflecting 300 shared\_initial\_generation calls plus a single repair invocation (stage\_counts: shared\_initial\_generation=300, repair=1). No deterministic exclusions for contamination were triggered by the preregistered checks; the execution and analysis artifacts enumerating per‑task normalized responses, validator traces, and repair traces are archived and are the primary source for the numerical counts below.

Primary, pair‑level outcomes (artifact counts): From the archived paired contingency artifact, the per‑condition success tallies are: baseline\_success\_count = 299/300 (99.6667\%), guarded\_success\_count = 300/300 (100.0\%). The full paired contingency table derived from the archived analysis artifact is n11 = 299 (both arms success), n10 = 1 (guarded success only), n01 = 0 (baseline success only), n00 = 0 (both fail). These contingency counts are reproduced verbatim in the archived artifact analysis/paired\_contingency\_table.csv and underpin all confirmatory inference reported here.

Confirmatory inferential results (preregistered analyses using archived artifacts): Applying the preregistered exact McNemar two‑sided test to the observed contingency yields p = 1.00. The observed absolute paired difference (guarded - baseline success rate) = 1/300 \ensuremath{\approx} 0.0033333. Paired bootstrap percentile 95\% CI for the absolute paired difference computed per the preregistered procedure (10,000 resamples, bootstrap\_seed = 1729) = [0.00, 0.01]. Repair outcomes: the single allowed automated repair corrected the sole invalid baseline output (repair success 1/1). The preregistered exact Clopper--Pearson 95\% interval for that conditional repair success is [0.025, 1.0], reflecting the wide uncertainty when only a single initial failure was observed.

Post‑hoc sensitivity and interpretive bounds grounded in observed discordant counts: The total number of discordant pairs observed is n10 + n01 = 1. Under the paired within‑task sampling semantics used, this discordant total imposes a strict upper bound on the resolvable absolute paired difference equal to (n10 - n01)/N = 1/300 \ensuremath{\approx} 0.00333. Practically, that bound means the empirical data cannot support detection of large absolute improvements (for example, the preregistered >=50\% relative reduction target, which was premised on an assumed baseline pass rate of 0.40 and corresponds to an absolute increase of \ensuremath{\approx}0.30) because achieving such an effect with N=300 would have required on the order of many tens of discordant pairs (the preregistration's heuristic estimate suggested \ensuremath{\approx}90 discordant pairs would be needed to observe such an effect size under conservative discordant assumptions). The archived contingency artifact and the bootstrap CI concretely illustrate this limitation: the 95\% CI for the absolute paired difference terminates at 0.01, well below the preregistered planning effect.

Why the exact McNemar p = 1.00 arises here (artifact‑grounded intuition): Exact McNemar inference bases significance on the imbalance of the two discordant cell counts (n10 versus n01). With only one discordant pair (1 vs 0) the exact two‑sided test computes a tail probability that includes that single outcome and any outcomes at least as extreme under the null; with such minimal discordance the resulting two‑sided p is not small and in this case equals 1.00 as reported by the preregistered executor. The archived paired\_contingency\_table.csv and analysis/analysis\_log.jsonl record the discordant counts and the exact test invocation that produced this p‑value.

Diagnostic traceability and representative artifact observations: For every episode the deterministic validator trace archived fields including normalized\_configuration, parsed\_command\_count, required\_command\_count, violation\_codes, and repair\_applied flags. The single baseline failure is fully traceable in the archived task artifact (execution/scoring/task-000XXX-baseline.json entry) and shows the validator‑reported violation that triggered the one repair invocation; the archived repair prompt/response pair documents the change applied and the validator's post‑repair canonicalized configuration which satisfied the full intent constraints. These per‑task artifacts are the authoritative evidence for the repair outcome and are enumerated in the execution manifest and representative task listings.

Implications of shared\_initial\_candidate pairing for observed information: The preregistered shared‑candidate design intentionally reused the same generated plan across both arms to isolate the validator/repair effect on a fixed candidate. While operationally motivated, this design choice increases within‑pair correlation and consequently reduces expected discordant counts compared with an independent‑draws design; fewer discordant pairs directly reduce the statistical information available to McNemar inference for a given N. The observed near‑ceiling baseline (299/300) in combination with the shared‑candidate pairing therefore explains the very small observed discordant total and the resulting inability to confirm the preregistered large relative effect. These relationships between pairing semantics, baseline prevalence, and discordant counts are visible in the archived deterministic\_reconciliation.json and analysis/condition\_summary.csv artifacts and informed the post‑hoc sensitivity statements above.

\section{Discussion}
We expand the Discussion to emphasize concrete, artifact‑grounded interpretation, diagnostic lessons from the run, and practical guidance for designing future confirmatory NetOps studies.

Quantitative interpretation tethered to archived counts. The archived contingency (n11=299, n10=1, n01=0, n00=0) directly constrains inferential outcomes: the paired difference equals (n10 - n01)/N = 1/300 \ensuremath{\approx} 0.00333, and any two‑sided exact McNemar test necessarily conditions on the observed discordant count (here 1). With only one discordant pair, the sampling distribution under the null contains insufficient extreme outcomes to produce a small p‑value; operationally this is why the exact McNemar test returns p = 1.00 for our observed pattern. The paired bootstrap CI [0.00, 0.01] (10,000 resamples, seed 1729) expresses the same limitation in interval form: the data are consistent only with extremely small absolute improvements given the observed discordance. These numerical facts are archived and reproduce directly from the execution manifest and analysis artifacts.

Why preregistered planning assumptions matter. Our preregistered power plan assumed a baseline pass rate = 0.40; that planning choice drove an expectation of many discordant pairs under a substantial relative reduction target. The actual baseline (0.9967) collapses that expectation. The practical lesson is that power planning for generate‑verify repair experiments must explicitly consider plausible baseline prevalences and the chosen pairing semantics (shared candidate vs independent draws), because both jointly determine expected discordant counts. A simple arithmetic relation is useful when translating target absolute paired differences into discordant‑pair targets: absolute paired difference = (n10 - n01)/N, so for N=300 an absolute increase of 0.30 requires (n10 - n01) \ensuremath{\approx} 90 net discordant pairs favoring guarded outputs. The archived artifacts provide the observed (near‑zero) discordant count and therefore an artifact‑grounded explanation of why the preregistered >=50\% relative reduction could not be detected.

Pairing semantics: operational rationale and inferential tradeoffs. We preregistered a shared\_initial\_candidate design to emulate a realistic operational workflow: operators often review and repair a single candidate plan rather than comparing independently drawn proposals. This choice improves ecological validity for the intended operational claim (effect of validator+repair on a single generated plan) but reduces statistical signal because within‑pair correlation lowers expected discordant counts. In contrast, an independent‑draws design yields larger expected discordance for a given marginal success probability and thus greater power for McNemar‑style paired tests; however, it addresses a different operational question (expected net benefit of switching from one generator draw to another) and increases heterogeneity. Future preregistrations should explicitly link the pairing semantics to the estimand of operational interest and incorporate that link into power and sample‑size computations.

Repair effectiveness as a conditional estimand and sample‑size implications. H2 framed repair success conditional on initial failure. The run observed repair success 1/1, but that conditional estimate is extremely imprecise (Clopper--Pearson 95\% CI = [0.025, 1.0]). Practically, when repair effectiveness is a primary operational claim, study designers should power the experiment for the conditional probability P(repair succeeds | initial failure) by ensuring a planned number of initial failures (not merely overall N). This can be achieved either by calibrating task difficulty to increase baseline failure prevalence or by targeted adversarial augmentation of tasks so that the expected number of initial failures meets the sample‑size requirements for the desired CI width.

Diagnostics enabled by deterministic validator traces. The frozen deterministic validator produced structured per‑task traces (normalized\_configuration, parsed\_command\_count, required\_command\_count, violation\_codes, repair\_applied). These traces are essential for post‑hoc failure‑mode analysis: they identify whether failures arise from syntax/canonicalization mismatches, missing commands, dependency‑order violations, or preserved‑state violations (e.g., touching protected interfaces). In our run the single baseline failure was diagnosed as a required‑sequence omission by the validator and repaired by applying the preregistered repair adapter; the repair trace shows the minimal edit needed to satisfy intent constraints. Maintaining such per‑task structured traces is a practical requirement for explainable NetOps guardrails and for reproducible post‑mortems.

Construct validity and the role of dynamic emulation. Our validator performs deterministic static checks and canonicalization rather than executing configurations in devices or emulators; this design increases safety by avoiding live changes but limits construct validity with respect to runtime semantic properties (control‑plane convergence, routing dynamics, packet‑level reachability under transient state). If operational claims require guarantees about runtime consequences, future studies should integrate emulation or simulation stages that exercise control‑plane and dataplane effects and record dynamic assertions (e.g., convergence time, transient reachability). The archived manifest documents that our results are strictly conditional on the validator's static semantics.

Design recommendations grounded in observed artifacts. Based on the run artifacts and contingency counts, we recommend three concrete design levers for future confirmatory NetOps evaluations:
1) Calibrate task difficulty or perform adversarial augmentation to target a nontrivial baseline failure prevalence or a predetermined discordant‑pair count aligned with the planned absolute/relative effect size. Use deterministic task generators with adjustable difficulty knobs and preregister the target discordant count. 2) Explicitly preregister pairing semantics and repair budgets and incorporate their expected effects on discordant counts into power calculations; if the operational question aims to evaluate remediation of a single plan, prefer shared‑candidate pairing but then set task difficulty to produce more initial failures. 3) When runtime semantics matter, add an emulation/simulation validation stage and preregister how emulator outcomes map to success/failure to avoid post‑hoc redefinition of the estimand. All of these recommendations are traceable to artifacts in the execution manifest (task difficulty fields, validator violation codes, and the paired contingency table) and are therefore actionable.

Threats to validity revisited with artifact examples. Internal validity: the shared\_candidate design and deterministic replacement rules are preregistered but reduce the number of informative discordant pairs; this is evident in the archived stage\_counts and provider\_call audit. Construct validity: static validator checks may miss dynamic failure modes; the validator traces explicitly show which checks were applied and which dynamic behaviors were not exercised. External validity: the synthetic task bank and constrained vendor templates are enumerated in the task manifest and limit generalization across vendor heterogeneity and production topologies. Conclusion validity: the paucity of discordant pairs (one) creates high sampling uncertainty for large preregistered effects; the paired bootstrap and exact McNemar results in the archived analysis artifacts document this limitation quantitatively.

Operational implications for deploying guarded pipelines. The run shows that a frozen deterministic validator plus a bounded repair loop can correct at least some invalid candidates without human intervention (the single observed correction). However, whether such a pipeline materially reduces first‑pass misconfiguration in production depends on the baseline error prevalence and the validator's coverage of real failure modes. When baseline error is already rare, the primary operational value of a validator+repair may be as a safety net that prevents the rare catastrophic mistake rather than as a mechanism to substantially increase throughput. Organizations should therefore align evaluation goals---error‑reduction vs catastrophic‑risk mitigation---with study design choices (task bank selection, pairing semantics, emulation inclusion) and with the operational acceptance thresholds that govern whether automated repair may be applied without human approval.

Reproducibility and archival practice. This run demonstrates the utility of a full artifact archive: per‑task inputs, normalized responses, validator traces, repair prompts/responses, the paired contingency table, and the analysis code together enable independent recomputation of reported counts and intervals. The execution manifest documents model\_calls\_used = 301 and repair\_stage invocations = 1; these accounting numbers are minimal but sufficient to audit the reported outcomes. Future confirmatory NetOps work should adopt similar archival granularity to permit external verification and subsequent meta‑analysis.

\section{Limitations}
\begin{itemize}
\item Synthetic task bank and constrained vendor‑template scope limit external validity to broad production environments.
\item Single LLM (gpt-5-mini) and a single validator implementation restrict generality across model families and validator designs.
\item The preregistered power plan assumed a baseline pass rate = 0.40; the observed baseline (299/300) contradicts that assumption and removed headroom to detect the preregistered >=50\% relative reduction.
\item Sampling semantics (shared\_initial\_candidate) reused the same initial candidate across paired arms, increasing within‑pair correlation and reducing discordant pairs available to McNemar inference.
\item Validation used deterministic configuration/constraint checks rather than full dynamic emulation of runtime semantic effects; actual runtime consequences of configurations were not executed and remain unmeasured.
\item H2 repair‑success evidence is based on a single initially invalid case (1/1 \ensuremath{\rightarrow} 100\%), yielding a wide Clopper--Pearson 95\% CI = [0.025, 1.0]; larger counts of initial failures are required to estimate repair effectiveness precisely.
\end{itemize}

\section{Conclusion}
We executed a preregistered paired experiment (N=300 pairs) comparing single‑shot LLM generation (gpt-5-mini) to a guarded pipeline consisting of a frozen deterministic validator and a single automated repair iteration. Baseline single‑shot generation yielded 299/300 valid outputs and the guarded pipeline yielded 300/300 valid outputs. The single automated repair corrected the lone invalid baseline output, but the observed near‑ceiling baseline contradicted the preregistered planning assumption (baseline pass rate = 0.40) and removed statistical headroom to detect the preregistered >=50\% relative reduction. Future confirmatory NetOps evaluations should (a) calibrate task difficulty to produce sufficient initial failures or target discordant counts, (b) preregister pairing semantics and repair budgets explicitly, and (c) include emulation to measure runtime semantic effects when operational deployment claims are intended. All raw outputs, validator traces, repair traces, the execution manifest, and analysis artifacts are archived and enumerated in the Disclosure Statement to permit reproduction of the reported counts and intervals.

\section*{Disclosure Statement}
This study was executed autonomously after the autonomy lock under the archived master prompt (provenance/master\_prompt.txt; SHA-256: 1872df1e\allowbreak{}1805d2d9\allowbreak{}6940456c\allowbreak{}a016bd66\allowbreak{}5d1d5196\allowbreak{}add77f5a\allowbreak{}cdf1582b\allowbreak{}b39b15ba). Human role: the Research Director selected the topic family (Generative AI and LLMs for NetOps) and approved the master prompt prior to the autonomy lock; no manual human editing of manuscript technical content, figures, tables, or references occurred after the lock. Preregistration: cand-2026-001-A-prereg-v1 (the preregistration record was created before observing final outcomes; preregistration SHA-256 is recorded in the execution manifest). Declared model and tooling: gpt-5-mini (declared\_model\_version), adapter family hosted\_netops\_gvr\_v1, frozen deterministic validator/repair adapter deterministic\_netops\_validator\_v5. Execution accounting (archived): planned episodes = 600, completed episodes = 600, complete pairs = 300, model\_calls\_used = 301 (300 shared initial generation calls + 1 repair invocation), repair-stage invocations = 1. The execution manifest and analysis artifact bundle enumerate all raw model responses, per‑task validator traces (normalized configurations, parsed\_command\_count, required\_command\_count, violation\_codes, repair\_applied flags), repair prompts/responses, execution logs, and analysis artifacts. The archived execution manifest and analysis bundle are the authoritative source for the numeric counts reported in this manuscript. The autonomous pipeline performed literature search, preregistration, execution, analysis, AI‑peer‑review style checks, and manuscript drafting autonomously after the autonomy lock in accordance with the CNSM Agentic AI Researcher track requirements. The preregistered planning assumption used in power calculations (planned baseline pass rate = 0.40) is explicitly reported here because it materially affects interpretation of the confirmatory result.

\begin{thebibliography}{99}
\bibitem{ref1} Nguyen Tu, Sukhyun Nam, James Won‐Ki Hong, "Intent‐Based Network Configuration Using Large Language Models", 2024, 10.1002/nem.2313
\bibitem{ref2} Oghenekeno Hilkiah Okula, ""The Era of AI Agents for Network Engineering": Intent-Aware Auto Remediation for Network Disruption or Failures Using AI Agents, Large Language Models (LLM) and Model Context Protocol (MCP) Mechanisms", 2026, 10.36227/techrxiv.177004277.71795055/v1
\bibitem{ref3} Kunal Dhanda, "MedRAG: Retrieval-Augmented Generation for Medical QA-Comparing Base and RAG-Augmented LLM Performance on Evidence from Peer-Reviewed Clinical Research", 2026, 10.2139/ssrn.6608518
\bibitem{ref4} Goutam Adwant, Manjari Srivastav, "Evidence Coverage Evaluator: A Comprehensive Framework for Assessing Grounding Quality in Retrieval-Augmented Generation Systems", 2026, 10.36227/techrxiv.176784505.56675782/v1
\bibitem{ref5} Ryan Beckett, Aarti Gupta, Ratul Mahajan, David Walker, "A General Approach to Network Configuration Verification", 2017, 10.1145/3098822.3098834
\bibitem{ref6} Apostol Vassilev, Alina Oprea, Alie Jean Fordyce, Hyrum S. Anderson, "Adversarial machine learning :", 2024, 10.6028/nist.ai.100-2e2023
\end{thebibliography}

\end{document}
